% 06-ket-qua.tex — Chương 6: Kết quả áp dụng

\chapter{Kết quả áp dụng}\label{ch:ketqua}

\begin{tomtat}[Phạm vi của chương]
Chương này tổng hợp những gì đã được xây dựng trong hệ thống tính đến thời điểm báo cáo. Hệ thống chưa vận hành chính thức nên chương chưa có số liệu đo chất lượng mô hình và kết quả kinh doanh; các số liệu này sẽ được bổ sung ở báo cáo sau.
\end{tomtat}

\section{Kết quả về sản phẩm}

\begin{enumerate}
  \item Dây chuyền khép kín từ dữ liệu thô đến danh sách khách hàng ưu tiên. Sáu luồng công việc tự động nối tiếp nhau, từ nạp dữ liệu đến chấm điểm, không cần thao tác tay giữa các bước.
  \item Danh sách khách hàng rủi ro cao theo chu kỳ. Mỗi lần chạy, hệ thống ghi bảng rủi ro vào cơ sở dữ liệu và xuất tệp CSV để đội chăm sóc dùng ngay; danh sách giới hạn ở nhóm 5\% rủi ro cao nhất, đủ ngắn để xử lý.
  \item Mỗi khách hàng có lý do rủi ro. Tối đa ba lý do, thuộc tám nhóm nghiệp vụ, kèm chỉ số hiện tại, mức nền và mức thay đổi.
  \item Cơ chế tự bảo vệ chất lượng. Mô hình được huấn luyện lại định kỳ (chu kỳ thiết kế khi đủ dữ liệu), chỉ thay thế khi tốt hơn và qua rào chắn an toàn; khi có sự cố, mô hình cũ vẫn tiếp tục phục vụ.
  \item Giám sát và truy vết đầy đủ. Nhật ký từng lần chạy, theo dõi độ lệch dữ liệu và điểm số, dọn dẹp tự động theo quý.
  \item Chạy trên hạ tầng sẵn có. Toàn bộ hệ thống đóng gói theo container, dùng PostgreSQL của đơn vị và có quy trình đóng gói, phát hành, hoàn tác phiên bản.
\end{enumerate}

\section{Kết quả về phương pháp}

{\tablefont
\begin{longtable}{@{}L{3.8cm}L{4.2cm}>{\raggedright\arraybackslash}p{7.21cm}@{}}
\caption{Điểm khác biệt về phương pháp so với cách làm thông thường}\label{tab:khacbiet}\\
\toprule
\hdrrow\hh{Khía cạnh} & \hh{Cách làm thông thường} & \hh{Cách làm của sáng kiến}\\
\midrule
\endfirsthead
\toprule
\hdrrow\hh{Khía cạnh} & \hh{Cách làm thông thường} & \hh{Cách làm của sáng kiến}\\
\midrule
\endhead
\bottomrule
\endfoot
Định nghĩa rời bỏ & Chỉ đếm khách ngừng hẳn & Gồm cả khách giảm mạnh số đơn, doanh thu, hoặc xu hướng đi xuống\\
\rowa Nhãn học & Chỉ dùng một nguồn & Nhãn theo bốn quy tắc nghiệp vụ dựa trên đơn hàng thật\\
Kiểm tra mô hình & Chia ngẫu nhiên & Chia theo thời gian, có khoảng đệm, giữ tháng gần nhất làm bài thi cuối kỳ\\
\rowa Thay mô hình & Thay khi có mô hình mới & Chỉ thay khi tốt hơn và qua rào chắn\\
Kết quả & Chỉ có điểm số & Điểm, thứ hạng và lý do bằng ngôn ngữ nghiệp vụ\\
\rowa Vận hành & Chạy tay từng bước & Tự động theo lịch, có nhật ký và giám sát\\
\end{longtable}}
\vspace{4pt}

\section{Số liệu đánh giá mô hình}

Mô hình đã được huấn luyện và kiểm tra theo thời gian trong quá trình phát triển, nhưng hệ thống chưa đưa vào vận hành chính thức và kết quả đo lường trên bài thi cuối kỳ chưa được tổng hợp. Các thước đo gồm điểm F1, độ chính xác cảnh báo, độ bao phủ, PR-AUC và ROC-AUC sẽ được báo cáo ở giai đoạn sau, kèm tháng dữ liệu và phiên bản mô hình tương ứng; chỉ kết quả trên bài thi cuối kỳ, không dùng kết quả của các vòng tinh chỉnh, được dùng để báo cáo.

\section{Đánh giá mức độ hoàn thiện}

{\tablefont
\begin{longtable}{@{}L{5.2cm}L{3.0cm}>{\raggedright\arraybackslash}p{7.01cm}@{}}
\caption{Mức độ hoàn thiện theo hạng mục}\label{tab:hoanthien}\\
\toprule
\hdrrow\hh{Hạng mục} & \hh{Trạng thái} & \hh{Ghi chú}\\
\midrule
\endfirsthead
\toprule
\hdrrow\hh{Hạng mục} & \hh{Trạng thái} & \hh{Ghi chú}\\
\midrule
\endhead
\bottomrule
\endfoot
Nạp và kiểm tra dữ liệu & \textcolor{okgreen}{\bfseries Hoàn thành} & Chạy hằng ngày\\
\rowa Tạo chỉ số, lọc điều kiện, gán nhãn & \textcolor{okgreen}{\bfseries Hoàn thành} & Nhãn theo quy tắc nghiệp vụ\\
Mô hình, tinh chỉnh, xét duyệt & \textcolor{okgreen}{\bfseries Hoàn thành} & Huấn luyện lại theo chu kỳ, thiết kế khi đủ dữ liệu\\
\rowa Chấm điểm và xuất danh sách & \textcolor{okgreen}{\bfseries Hoàn thành} & Bảng cơ sở dữ liệu + CSV\\
Giám sát và dọn dẹp & \textcolor{okgreen}{\bfseries Hoàn thành} & Độ lệch dữ liệu, độ lệch điểm số\\
\rowa Phân loại nguyên nhân chủ động/phản ứng & \textcolor{amber}{\bfseries Một phần} & Chỉ số đã có, chưa lưu nhãn riêng\\
Đổi tên cột tỷ lệ rời bỏ cho rõ nghĩa & \textcolor{amber}{\bfseries Một phần} & Đã ghi chú, đề xuất đổi tên\\
\rowa Xử lý khách mùa vụ, khách mới, khách phục hồi & \textcolor{warnred}{\bfseries Chưa làm} & Trong lộ trình (Chương~\ref{ch:morong})\\
Quản lý vòng đời mô hình & \textcolor{warnred}{\bfseries Chưa làm} & Trong lộ trình\\
\rowa Thử nghiệm đối chứng & \textcolor{warnred}{\bfseries Chưa làm} & Đề xuất ở Chương~\ref{ch:hieuqua}\\
\end{longtable}}
\vspace{4pt}
